\section{Main Theorem}\label{sec:main-theorem}

\begin{theorem}[Rounded response-tree representation]
\label{thm:main}
Under Assumptions~\ref{assump:sq-parameter-regime}
and~\ref{assump:universal-adversarial-sq}, let
\(K=\lceil1/\tau\rceil\),
\(G=\{-1+2j/K:0\le j\le K\}\), and
\(N=(K+1)^m\), and let \(\mathsf P_A\) be the pre-instance response-tree
law defined in \eqref{eq:response-tree-map-law}.  Then the same law works
simultaneously for every distribution \(\mathcal D\) on \(\mathcal X\)
and every \(h\in\mathcal H\), and
\[
\begin{aligned}
&\mathbb E_{\Phi\sim\mathsf P_A}
\left[
\inf_{w\in\mathbb R^N}
L_{\mathrm{tie}}\bigl(\mathcal D,h,\langle w,\Phi\rangle\bigr)
\right]\\
&\quad\le
\mathbb E_{R\sim\mu}
L_{\mathrm{tie}}(\mathcal D,h,g_{R,z^{\mathcal D,h,R}})\\
&\quad=
\mathbb E_{R\sim\mu}
L_{\mathrm{bin}}\bigl(
\mathcal D,h,
A_R^{\mathcal O^\rho_{\mathcal D,h}}(\mathcal D,h)
\bigr)
\le\varepsilon.
\end{aligned}
\tag{11}\label{eq:main-risk-chain}
\]
Consequently,
\[
\operatorname{dc}^{\mathrm{tie}}_\varepsilon(\mathcal H)
\le N
=\bigl(\lceil1/\tau\rceil+1\bigr)^m.
\tag{12}\label{eq:main-dimension-bound}
\]

Both bounds are exact and have no hidden constants.  The exposed variables
are \(m\), \(\tau\), \(\varepsilon\), \(K\), and \(N\); no dependence on
\(\mathcal X\), \(|\mathcal X|\), \(\mathcal H\), \(\mathcal D\),
\(h\), a tape, or a transcript is hidden.  The domain, class, learner,
tape law, rounding rule, coordinate order, and padding convention are fixed
before the instance.  The outer probability mode is expectation only over
\(R\sim\mu\), equivalently \(\Phi\sim\mathsf P_A\); the horizon mode is a
fixed finite budget of at most \(m\) adaptive queries; and the loss mode is
tie-penalized \(0\)-\(1\) risk, transferred exactly to binary \(0\)-\(1\)
risk on the selected coordinate.  In particular, \(m=0\) gives \(N=1\),
\(\varepsilon=0\) gives zero expected tie loss without a rounding
remainder, and \(\tau\ge1\) gives \(N=2^m\).
\end{theorem}
